This Master's Thesis evaluates whether \textit{Differential Machine Learning} (DML) improves, relative to a conventional MLP, the valuation and especially the Delta estimation of a fixed-strike arithmetic Asian option written on compute capacity. A daily NVIDIA H100 rental-price series is used as the main empirical reference, complemented by descriptive evidence from the Chinese market. Physical price dynamics are represented by a Log-Ornstein--Uhlenbeck process, while valuation is performed under a modelled pricing law. Reference prices and sensitivities are generated using Randomized Quasi-Monte Carlo and a pathwise Delta estimator. MLP and DML share the same architecture, samples and training protocol, and are compared across four training-set sizes and five paired seeds.

With $65\,536$ training states, DML reduces test price RMSE from $0.001624$ to $0.001516$, although the MLP achieves a slightly lower MAE. The difference is substantially larger for Delta: RMSE falls from $0.008528$ to $0.001817$ and MAE from $0.001234$ to $0.000216$. This improvement carries over to local hedging: for a $1\%$ spot shock, DML attains an RMSE of $1.058\times10^{-6}$ versus $6.688\times10^{-6}$ for the MLP and $1.355\times10^{-4}$ without hedging. In the dynamic hedge using model-implied compute forwards, DML strongly improves on the MLP, but a small number of extreme paths prevents it from outperforming the unhedged position in aggregate RMSE.

The results show that differential supervision robustly improves sensitivity learning and its local risk-management value, while also demonstrating that a more accurate Greek is not by itself sufficient to guarantee superior dynamic hedging performance.

\vspace{0.5em}
\noindent\textbf{Keywords:} Differential Machine Learning; Asian options; compute capacity; Randomized Quasi-Monte Carlo; Delta; hedging; Log-OU.
